\documentclass[10pt,a4paper]{amsart}
\usepackage[margin=19mm]{geometry}
\usepackage{amsmath,amssymb,mathtools}
\usepackage[scr=rsfs,cal=euler]{mathalfa}
\usepackage{xfrac}
\usepackage[unicode,hidelinks,bookmarks=false]{hyperref}
\newcommand{\ie}{i.e.\ }
\newcommand{\cf}{cf.\ }
\newcommand{\resp}{respectively\ }
\newcommand{\Subsection}{\subsection}
\providecommand{\nobreakdash}{\nobreak\hbox{-}}
\setlength{\emergencystretch}{2em}
\multlinegap0pt
\usepackage{bm}
\usepackage{bbm}
\usepackage{dsfont}
\usepackage{xspace}
\usepackage{cancel}
\usepackage{mathtools}
\usepackage{amsthm}
\makeatletter
\@ifundefined{theorem}{\newtheorem{theorem}{Theorem}}{}
\@ifundefined{lemma}{\newtheorem{lemma}[theorem]{Lemma}}{}
\@ifundefined{prop}{\newtheorem{prop}[theorem]{Proposition}}{}
\makeatother
\title[Uniqueness in the Plateau problem near Quadratic cones]{Uniqueness in the Plateau problem near Quadratic cones}
\author{Vishnu Nandakumaran, G\'abor Sz\'ekelyhidi}
\date{Source: arXiv:2509.15503v1 (CC BY 4.0)}
\begin{document}
\maketitle

\noindent\emph{Source and licence.} Uniqueness in the Plateau problem near Quadratic cones. Version arXiv:2509.15503v1 (CC BY 4.0), distributed under the Creative Commons Attribution 4.0 International licence, as recorded in the arXiv licence metadata for this exact version and shown on the abstract page https://arxiv.org/abs/2509.15503v1. Full rights notice: licenses/arxiv-2509.15503-CC-BY-4.0.txt.\par\smallskip

\noindent\textit{Editorial scope.} Counted unit: the Lemma whose source label is \texttt{lemma:k\_i bounded} in the article (source0.tex lines 403--406, source label \texttt{lemma:k\_i bounded}), delivered together with the article's own complete original proof of it (source0.tex lines 407--431, one \texttt{proof} environment of 25 lines). No mathematics is written, repaired or re-derived: statement and proof are the article's own text line for line. Required context retained uncounted: prelim:GammaBound (lines 247--257); prop:FromAllard (lines 278--281). Every definition, display and label of those lines is retained unchanged. Omitted from this package and not redistributed: 1--246, 258--277, 282--402, 432--525. Nothing inside a retained excerpt is omitted or altered: every retained body is delivered exactly as the article's lines are written, so no omission receipt of any kind accompanies this package. Citations: the retained excerpts carry no citation command at all, so no bibliography entry of the article is transcribed and no reference is redistributed. Reference adaptation: none was needed and none was made; every reference inside the delivered unit resolves inside it, so no unresolved reference or citation is produced. Printed slips kept literal: no printed word, symbol, hypothesis or number of the counted statement or of its proof was altered. Layout and class: the article's own class is \texttt{amsart} with packages including amsmath, amssymb, amsthm, verbatim, color, babel, amsfonts, graphicx, mathtools, float, xcolor, biblatex, hyperref; the wrapper is the lane's pinned \texttt{amsart} editorial wrapper at 10pt on A4, which restates the theorem-like environments the retained text needs and adds the article's own macro definitions that the retained text needs (none), with extra wrapper packages (bm, bbm, dsfont, xspace, cancel, mathtools). The delivered numbering is the unit's own. Assets: no retained span contains a figure environment, a graphics-inclusion command, a PDF-page inclusion, a table or a TikZ picture; no image file is copied into this package, the package declares no asset dependency and no image file is redistributed with it. Licence: version arXiv:2509.15503v1 of this article is distributed under the Creative Commons Attribution 4.0 International licence, as recorded in the arXiv licence metadata for this exact version and shown on the abstract page https://arxiv.org/abs/2509.15503v1. A full rights notice is delivered at licenses/arxiv-2509.15503-CC-BY-4.0.txt. Characters: every retained excerpt is pure ASCII, so the delivered engine, XeLaTeX with its default Latin Modern text font, reports no missing character.
\begin{lemma} \label{prelim:GammaBound}
We have the following, 
\begin{enumerate}
    \item Let $w$ be a Jacobi field on $S_\lambda$ for $\lambda\in \mathbb R\setminus\{0\}$ such that 
$$ |w(x)|<C|x|^{\frac{2-n}{2}}$$
for $|x| > 1$.  Then $w=0$.

\item Let $w$ be a Jacobi field on $\mathcal C_1 = \mathcal{C} \cap B_1$ with $w= 0$ on $\Sigma$ and $|w(x)|<C|x|^{\frac{2-n}{2}}$ for $|x| < 1$. Then $w=0$  on $\mathcal C_1$.
\end{enumerate}

\end{lemma}

\begin{prop}\label{prop:FromAllard}
    Let $c_0, r_0 > 0$ be small. There exists an $\epsilon_0 > 0$ depending on $c_0, r_0$, with the following property. Let $\Vert g\Vert_{C^{2,\alpha}(\Sigma)} < \epsilon_0$ and suppose that $N$ is a stationary integral varifold in $B_1$ with boundary $\Sigma'$ given by the spherical graph of $g$ over $\Sigma$. In addition suppose that $\mu_N(B_1) < \frac{3}{2} \mu_\mathcal C (B_1)$. 
    Then in the annulus $B_1\setminus B_{r_0}$ the varifold $N$ is the spherical graph of a function $v$ over $\mathcal{C}$ with $\Vert v\Vert_{C^{2,\alpha}} < c_0$. 
\end{prop}

\begin{lemma}\label{lemma:k_i bounded}
    If the $M^i_1, M^i_2$ are distinct, then 
    there is a constant $K > 0$ such that $k_i < K$ for all sufficiently large $i$. 
\end{lemma}

\begin{proof}
We argue by contradiction. Suppose that $k_i\rightarrow \infty$ ($\rho_0^{k_i}\rightarrow 0$) or $k_i=\infty$ along a subsequence. In either scenario we have a sequence of positive integers $l_i\rightarrow \infty$ such that 
For all $0\leq k<l_i$, 
\begin{equation}\label{eqn: M_j distinct}
 \|u_i\|_{\rho_0, k+1}< \rho_0^{\frac{2-n}{2}}\|u_i\|_{\rho_0, k}<\rho_0^{k(\frac{2-n}{2})}\|u_i\|_{\rho_0, 1}.
\end{equation}
We assume $u_i\neq 0$ and define $u_i^o:\hat M_1^i\cap B_1(x_1^i)\backslash B_{\rho_0^{k_i}}\rightarrow \mathbb R$ as
$$ u_i^o\coloneqq\frac{u_i}{\|u_i\|_{L^2(\hat M_1^i\cap B_1(x_1^i)\backslash B_{\rho_0^2})}}. $$

We claim that there exists a subsequence such that $u_i^o \rightarrow w$, where $w$ is a Jacobi field on $\mathcal C \cap B_1$, with $w=0$ on $\Sigma$ and $|w(x)| \leq C |x|^{\frac{2-n}{2}}$. Here the convergence is $C^{2,\alpha}$ on compact subsets of $\mathcal{C}\cap (B_1\setminus \{0\})$, up to the boundary $\partial B_1$, and in fact smooth on the interior. In order to show this, we will show that for large $i$ the $u^o_i$ satisfy uniform $C^{k,\alpha}$ estimates on any annulus $B_1(x^i_1)\setminus B_r$. Recall that from Proposition~\ref{prop:FromAllard} we know that as $i\to \infty$, the $\hat{M}^i_j$ restricted to any such annulus converge smoothly to $\mathcal{C}\cap (B_1\setminus B_r)$. 

To see this, note that
\[ \Vert u_i\Vert_{\rho_0,1} \leq C(\rho_0) \Vert u_i\Vert_{L^2(\hat{M}^i_1\cap B_1(x^i_1))\setminus B_{\rho_0^2}}. \]
From \eqref{eqn: M_j distinct} we then find that
\begin{equation}\label{eq:L2bound} \Vert u^o_i\Vert_{\rho_0, k+1} < C(\rho_0) \rho_0^{k(2-n)/2}, 
\end{equation}
for $k < l_i$. 

On any fixed annulus $B_1(x^i_1)\setminus B_r$ we can write the equation for $u_i$ as $\mathcal{M}_{\hat{M}_1^i}(u_i)=0$, and $u_i=0$ on $\partial B_1(x^i_1)$. In other words
\begin{equation}\label{eq:M=0} \mathcal{L}_{\hat{M}^i_1}(u_i) + Q_{\hat{M}^i_1}(u_i) = 0. 
\end{equation}
Since on the annulus $\Vert u_i\Vert_{C^{2,\alpha}}\to 0$ as $i\to\infty$, and $Q$ contains quadratic or higher order terms, it follows that the equation can be viewed as a homogeneous linear equation for $u_i^o$ with $C^\alpha$ coefficients (which depend on $u_i$). The interior and boundary Schauder estimates, together with the $L^2$ bounds \eqref{eq:L2bound}, imply that the $u_i^o$ satisfy uniform $C^{k,\alpha}$ bounds up to the boundary on any annulus $B_1(x^i_1)\setminus B_r$. In particular, up to choosing a subsequence, and letting $r\to 0$, we can extract a limit $w = \lim_{i\to\infty} u^o_i$ on $\mathcal{C}\cap B_1\setminus \{0\}$. Letting $i\to \infty$ in Equation \eqref{eq:M=0} implies that  $\mathcal{L}_{\mathcal{C}}(w)=0$, while the convergence along the boundary implies $w=0$ on $\Sigma$. 

Using the interior estimates, from \eqref{eq:L2bound} we can deduce that  $|w(x)|\leq C|x|^{\frac{2-n}{2}}.$ By Lemma~\ref{prelim:GammaBound} this leads to  $w\equiv 0$ which is a contradiction. 
\end{proof}

\end{document}
